\newpage
\appendix

\section{Python验证代码}

以下为用于数据拟合和可视化的Python代码

\subsection*{幂律拟合分析代码 (fit\_analysis.py)}

\inputpython{codes/metabolic_rate_codes/fit_analysis.py}{1}{40}

\inputpython{codes/metabolic_rate_codes/fit_analysis.py}{113}{151}

\subsection*{可视化分析代码 (visualization.py)}

\inputpython{codes/metabolic_rate_codes/visualization.py}{1}{50}

\inputpython{codes/metabolic_rate_codes/visualization.py}{87}{100}

\subsection*{交叉验证分析代码 (cross\_validation.py)}

\inputpython{codes/metabolic_rate_codes/cross_validation.py}{1}{35}

\inputpython{codes/metabolic_rate_codes/cross_validation.py}{38}{70}

\section{数据来源说明}

\subsection*{静息心率数据}
来源于MSD Veterinary Manual (Merck Sharp \& Dohme Corp.)，该手册是全球兽医领域的权威参考文献。数据链接：https://www.msdvetmanual.com/multimedia/table/resting-heart-rates

\subsection*{生理参数数据}
来源于哈佛医学院BioNumbers项目，该项目汇集了各类生物学参数的标准值。数据包括小鼠、大鼠、兔、猫、狗、人等动物的基础生理参数。

\subsection*{基础代谢率数据}
部分来源于FAO/WHO联合发布的《Human Energy Requirements》报告，该报告提供了基于体重和年龄预测BMR的标准公式。

\section{源代码}
本课程的相关作业和结课报告的源代码已上传至GitHub仓库，链接为：\url{https://github.com/Eitherwoods/2026SEUBSMA_ML-BME.git}
